%% ============================================================
%% preamble.tex — ARES Mission Workbook
%% Editorial Flat 스타일 (07_ResearchNote 양식 준수)
%% XeLaTeX 컴파일: xelatex main.tex
%% ============================================================

%% ----- 한국어 지원 -----
\usepackage{kotex}

%% ----- 폰트 (XeLaTeX, KoPub 패밀리) -----
%% \userfontpath를 사용 환경에 맞게 수정할 것.
%%   macOS:   /Users/<USER>/Library/Fonts/
%%   Windows: C:/Users/<USER>/AppData/Local/Microsoft/Windows/Fonts/
%%            또는 C:/Windows/Fonts/
\usepackage{fontspec}
\defaultfontfeatures{Mapping=}
\newcommand{\userfontpath}{C:/Users/young/AppData/Local/Microsoft/Windows/Fonts/}

\setmainfont{KoPub Batang}[
  Path=\userfontpath,
  Scale=0.92, Mapping=, LetterSpace=-2.0,
  UprightFont = KoPub Batang Medium.ttf,
  BoldFont    = KoPub Batang Bold.ttf,
  ItalicFont  = KoPub Batang Medium.ttf,
  ItalicFeatures = {FakeSlant=0.15},
  BoldItalicFont = KoPub Batang Bold.ttf,
  BoldItalicFeatures = {FakeSlant=0.15}
]
\setsansfont{KoPub Dotum}[
  Path=\userfontpath,
  Scale=0.92, Mapping=, LetterSpace=-2.0,
  UprightFont = KoPub Dotum Medium.ttf,
  BoldFont    = KoPub Dotum Bold.ttf
]
\setmonofont{D2Coding}[
  Path=\userfontpath,
  Scale=0.88,
  Extension = .ttc,
  UprightFont = D2Coding-Ver1.3.2-20180524-all
]
\setmainhangulfont{KoPub Batang}[
  Path=\userfontpath,
  Scale=0.92, Mapping=, LetterSpace=-2.0,
  UprightFont = KoPub Batang Medium.ttf,
  BoldFont    = KoPub Batang Bold.ttf
]
\setsanshangulfont{KoPub Dotum}[
  Path=\userfontpath,
  Scale=0.92, Mapping=, LetterSpace=-2.0,
  UprightFont = KoPub Dotum Medium.ttf,
  BoldFont    = KoPub Dotum Bold.ttf
]
\newhangulfontfamily\hanguldotum{KoPub Dotum}[
  Path=\userfontpath,
  Scale=0.92, LetterSpace=-2.0,
  UprightFont = KoPub Dotum Medium.ttf,
  BoldFont    = KoPub Dotum Bold.ttf
]
\setmonohangulfont{KoPub Batang}[
  Path=\userfontpath,
  Scale=0.92, LetterSpace=-2.0,
  UprightFont = KoPub Batang Medium.ttf,
  BoldFont    = KoPub Batang Bold.ttf
]

%% ----- 유럽 특수문자 폴백 -----
\usepackage{ucharclasses}
\newfontfamily\latinfont{Times New Roman}[Scale=0.92]
\makeatletter
\newcommand{\prevfontfam}{}
\newcommand{\saveandlatinfont}{\edef\prevfontfam{\f@family}\latinfont}
\newcommand{\restoreprevfont}{\fontfamily{\prevfontfam}\selectfont}
\makeatother
\setTransitionsFor{LatinSupplement}{\saveandlatinfont}{\restoreprevfont}
\setTransitionsFor{LatinExtendedA}{\saveandlatinfont}{\restoreprevfont}
\setTransitionsFor{LatinExtendedB}{\saveandlatinfont}{\restoreprevfont}

\newcommand{\dash}{\kern1pt{\latinfont\char"2014}\kern1pt}
\newcommand{\ndash}{{\latinfont\char"2013}}
\usepackage{newunicodechar}
\newunicodechar{–}{{\latinfont\char"2013}}
\newunicodechar{—}{{\latinfont\char"2014}}

%% ----- 목차 깊이 -----
\setcounter{tocdepth}{2}
\setcounter{secnumdepth}{1}  %% subsection(미션) 번호는 본문 텍스트 ``미션 N''으로 충분 → 카운터 비표시

%% ----- 줄간격/문단 -----
\linespread{1.12}
\setlength{\parskip}{0.85em}
\setlength{\parindent}{0pt}
\raggedbottom

%% ----- 페이지 레이아웃 -----
\usepackage{geometry}
\geometry{
  paperwidth=190mm, paperheight=240mm,
  top=24mm, bottom=28mm,
  inner=24mm, outer=22mm,
  headheight=14pt, headsep=8mm,
  footskip=28pt
}

%% ----- 수식 -----
\usepackage{amsmath, amssymb, amsthm}
\usepackage{mathtools}
\theoremstyle{definition}
\newtheorem{definition}{정의}[section]
\newtheorem{theorem}{정리}[section]
\newtheorem{lemma}{보조정리}[section]
\newtheorem{proposition}{명제}[section]
\newtheorem{example}{예제}[section]
\theoremstyle{remark}
\newtheorem*{remark}{Remark}

%% ----- 그래픽 -----
\usepackage{graphicx}
\usepackage{tikz}
\usepackage{pgfplots}
\pgfplotsset{compat=1.18}
\usetikzlibrary{arrows.meta, positioning, shapes.geometric, calc, shadows, fit, backgrounds}

%% ----- 표 -----
\usepackage{booktabs}
\usepackage{tabularx}
\usepackage{multirow}
\usepackage{array}

%% ----- 색상 정의 (Editorial Flat 팔레트) -----
\usepackage{xcolor}
\definecolor{deepblue}{RGB}{22, 56, 100}
\definecolor{skyblue}{RGB}{70, 130, 200}
\definecolor{accentcopper}{RGB}{200, 105, 50}
\definecolor{warmgold}{RGB}{218, 170, 65}
\definecolor{rulegray}{RGB}{210, 214, 220}
\definecolor{inkgray}{RGB}{55, 65, 80}
\definecolor{codebackground}{RGB}{249, 250, 252}
\definecolor{neutralgray}{RGB}{95, 105, 120}

%% 박스 색상
\definecolor{keypointmain}{RGB}{190, 145, 35}
\definecolor{cautioncolor}{RGB}{195, 85, 70}
\definecolor{labcolor}{RGB}{40, 140, 125}
\definecolor{notecolor}{RGB}{105, 115, 175}

%% ----- 코드 (minted, Pygments 기반) -----
%% listings는 UTF-8 한글을 바이트 단위로 처리하여 한글-숫자 경계에서
%% 글리프 박스 순서가 뒤집히는 알려진 문제가 있다. minted는 코드를
%% Pygments(Python)가 LaTeX 바깥에서 토큰화 후 결과만 넘기므로 한글이
%% 안전하게 처리된다.
%% 컴파일: xelatex -shell-escape main.tex
\usepackage{minted}
\usemintedstyle{friendly}
\setminted{
  fontsize=\footnotesize,
  bgcolor=codebackground,
  frame=leftline,
  framerule=0.8pt,
  rulecolor=\color{accentcopper},
  breaklines=true,
  breakanywhere=true,
  tabsize=4,
  linenos=false,
  autogobble=true,
}
\setminted[python]{
  python3=true,
}

%% ----- 박스 환경 (tcolorbox) -----
\usepackage{tcolorbox}
\tcbuselibrary{skins, breakable, listings}

%% 칩(chip) 헬퍼
\newcommand{\boxchip}[2]{%
  \tikz[baseline=(c.base)]{%
    \node[fill=#1, text=white, inner xsep=6pt, inner ysep=2pt,
          rounded corners=0pt, sharp corners]
      (c) {\sffamily\bfseries\fontsize{8}{9}\selectfont #2};%
  }\ignorespaces%
}
\newcommand{\boxchipwithtitle}[3]{%
  \boxchip{#1}{#2}\enspace{\sffamily\bfseries\footnotesize\color{deepblue}#3}%
}

%% 요점 박스 (카퍼)
\newtcolorbox{keypoint}{
  enhanced, breakable,
  sharp corners, arc=0pt,
  frame hidden, boxrule=0pt,
  colback=accentcopper!6,
  borderline west={2pt}{0pt}{accentcopper},
  before upper={\boxchip{accentcopper}{요점}\par\vspace{4pt}},
  top=6pt, bottom=6pt, left=12pt, right=10pt,
  before skip=1em, after skip=0.9em,
}

%% 주의 박스 (코럴)
\newtcolorbox{caution}{
  enhanced, breakable,
  sharp corners, arc=0pt,
  frame hidden, boxrule=0pt,
  colback=cautioncolor!6,
  borderline west={2pt}{0pt}{cautioncolor},
  before upper={\boxchip{cautioncolor}{주의}\par\vspace{4pt}},
  top=6pt, bottom=6pt, left=12pt, right=10pt,
  before skip=1em, after skip=0.9em,
}

%% 잠깐 박스 (바이올렛)
\newtcolorbox{notebox}[1]{
  enhanced, breakable,
  sharp corners, arc=0pt,
  frame hidden, boxrule=0pt,
  colback=notecolor!6,
  borderline west={2pt}{0pt}{notecolor},
  before upper={\boxchipwithtitle{notecolor}{잠깐}{#1}\par\vspace{4pt}},
  top=6pt, bottom=6pt, left=12pt, right=10pt,
  before skip=1em, after skip=0.9em,
}

%% 실습/실험 박스 (티)
\newtcolorbox{labbox}[1]{
  enhanced, breakable,
  sharp corners, arc=0pt,
  frame hidden, boxrule=0pt,
  colback=labcolor!6,
  colbacktitle=labcolor!22,
  coltitle=deepblue,
  fonttitle=\bfseries\sffamily\large,
  title={#1},
  toptitle=8pt, bottomtitle=8pt, lefttitle=14pt, righttitle=12pt,
  top=12pt, bottom=12pt, left=14pt, right=12pt,
  before skip=1em, after skip=0.9em,
}

%% 인용/강조 박스 (카퍼 헤어라인)
\newtcolorbox{bodyquote}{
  enhanced,
  frame hidden,
  borderline west={2pt}{0pt}{accentcopper},
  colback=white,
  arc=0pt, sharp corners,
  left=12pt, right=6pt, top=4pt, bottom=4pt,
  enlarge left by=6mm, enlarge right by=6mm,
  width=\linewidth-12mm,
  before skip=0.7em, after skip=0.7em,
  fontupper=\small\itshape\color{inkgray}
}

%% 절 요약 박스
\newtcolorbox{summary}[1][요약]{
  enhanced jigsaw, breakable,
  arc=2mm,
  colback=black!4,
  colframe=black!55,
  boxrule=0.6pt,
  title={#1},
  colbacktitle=black!55,
  coltitle=white,
  fonttitle=\sffamily\bfseries,
  top=6pt, bottom=6pt, left=12pt, right=10pt,
  before skip=1em, after skip=0.9em,
}

%% 스토리(대화) 박스 — 미션 도입부용
\newtcolorbox{storybox}{
  enhanced, breakable,
  sharp corners, arc=0pt,
  frame hidden, boxrule=0pt,
  colback=skyblue!6,
  borderline west={2pt}{0pt}{skyblue},
  before upper={\boxchip{skyblue}{스토리}\par\vspace{4pt}},
  top=6pt, bottom=6pt, left=12pt, right=10pt,
  before skip=1em, after skip=0.9em,
}

\usepackage{etoolbox}
\BeforeBeginEnvironment{keypoint}{\par}
\BeforeBeginEnvironment{summary}{\par}
\BeforeBeginEnvironment{caution}{\par}
\BeforeBeginEnvironment{bodyquote}{\par}
\BeforeBeginEnvironment{storybox}{\par}

%% (minted는 linenos가 기본 false이므로 박스 내부 행번호 토글 불필요)

%% ----- 헤더 / 푸터 -----
\usepackage{fancyhdr}
\newcommand{\evenfootblock}{%
  \begin{tabular}[b]{@{}l@{}}
    {\color{accentcopper}\rule{7mm}{1.2pt}}\\[2pt]%
    \small\sffamily\thepage\enspace{\color{deepblue}\sffamily\bfseries\small ARES Mission Workbook}
  \end{tabular}%
}
\newcommand{\oddfootblock}{%
  \begin{tabular}[b]{@{}r@{}}
    {\color{accentcopper}\rule{7mm}{1.2pt}}\\[2pt]%
    {\color{deepblue}\sffamily\bfseries\small\leftmark}\enspace\small\sffamily\thepage
  \end{tabular}%
}
\renewcommand{\sectionmark}[1]{\markboth{\thesection\enspace #1}{}}
\renewcommand{\subsectionmark}[1]{}

\pagestyle{fancy}
\fancyhf{}
\renewcommand{\headrulewidth}{0pt}
\renewcommand{\footrulewidth}{0pt}
\fancyhfoffset[LE,RO]{10mm}
\fancyfoot[LE]{\evenfootblock}
\fancyfoot[RO]{\oddfootblock}

\fancypagestyle{plain}{%
  \fancyhf{}
  \renewcommand{\headrulewidth}{0pt}
  \renewcommand{\footrulewidth}{0pt}
  \fancyhfoffset[LE,RO]{10mm}
  \fancyfoot[LE]{\evenfootblock}
  \fancyfoot[RO]{\oddfootblock}
}

%% ----- 하이퍼링크 -----
\usepackage{hyperref}
\hypersetup{
  colorlinks=true,
  linkcolor=deepblue,
  citecolor=deepblue,
  urlcolor=deepblue,
  pdftitle={ARES Mission Workbook},
  pdfauthor={(주)코리아사이언스, 동명대학교 게임학부},
  pdfsubject={ARES 자율 로버 탐사 시스템 미션 워크북},
  pdfkeywords={ARES, 블록코딩, Blockly, 라즈베리파이 피코, 로봇, 교육}
}

%% ----- 목차 스타일 (tocloft) -----
\usepackage[titles]{tocloft}
\renewcommand{\cftsecfont}{\sffamily\bfseries\color{deepblue!90}}
\renewcommand{\cftsecpagefont}{\sffamily\bfseries\color{deepblue!90}}
\renewcommand{\cftsubsecfont}{\small\rmfamily\color{black!75}}
\renewcommand{\cftsubsecpagefont}{\small\rmfamily\color{black!75}}
\setlength{\cftbeforesecskip}{4pt}

%% ----- 섹션 제목 — 대형 카퍼 번호 + 네이비 타이틀 + 헤어라인 -----
\usepackage{titlesec}

\newcommand{\sectionboxfmt}[1]{%
  \noindent
  \begin{tikzpicture}[baseline=(t.base)]
    \node[inner sep=0pt, outer sep=0pt] (n) at (0,0)
      {\sffamily\bfseries\fontsize{28}{28}\selectfont\color{accentcopper}\thesection};
    \node[anchor=base west, inner sep=0pt, outer sep=0pt] (t)
      at ([xshift=0.5em, yshift=2pt]n.base east)
      {\sffamily\bfseries\fontsize{18}{20}\selectfont\color{deepblue}#1};
  \end{tikzpicture}%
}
\titleformat{\section}[block]
  {\normalfont}
  {}{0pt}{\sectionboxfmt}
  [\vspace{4pt}{\color{rulegray}\hrule height 0.5pt}\vspace{6pt}]
\titlespacing*{\section}{0pt}{18pt plus 2pt}{10pt}

%% 별표 섹션(\section*, 예: \tableofcontents의 ``목차''): 번호 없이 제목만
\titleformat*{\section}
  {\sffamily\bfseries\fontsize{18}{20}\selectfont\color{deepblue}}

%% 차시(\section) 번호는 로마숫자, 미션(\subsection)은 순수 아라비아
\renewcommand{\thesection}{\Roman{section}}
\renewcommand{\thesubsection}{\arabic{subsection}}

%% subsection — 카퍼 텍스트 + 옅은 카퍼 라인
\titleformat{\subsection}
  {\normalfont\sffamily\large\bfseries\color{accentcopper}}
  {\thesubsection}{1em}{}
  [\vspace{1pt}{\color{accentcopper!35}\hrule height 0.4pt}]

\titleformat{\subsubsection}
  {\normalfont\sffamily\normalsize\bfseries\color{accentcopper}}
  {}{0pt}{}
  [\vspace{1pt}{\color{accentcopper!35}\hrule height 0.4pt}]

%% ----- 캡션 -----
\usepackage{caption}
\captionsetup{font={small,sf}, labelfont={small,sf,bf}}
\captionsetup[table]{position=above, skip=4pt}

%% ----- 유틸리티 -----
\usepackage{enumitem}
\usepackage{float}
\usepackage{placeins}
\usepackage{microtype}

%% ----- 취소선 -----
\usepackage[normalem]{ulem}
\newcommand{\st}[1]{\sout{#1}}

%% ----- 한국어 캡션 -----
\renewcommand{\contentsname}{목차}
\renewcommand{\figurename}{그림}
\renewcommand{\tablename}{표}
\renewcommand{\lstlistingname}{코드}

%% ----- 사용자 매크로 -----
\newcommand{\memref}[1]{\textsf{\small\color{neutralgray}[\detokenize{#1}]}}

%% 대화 라인 헬퍼 (\sayares{대사}, \sayalbi{대사})
\newcommand{\sayares}[1]{\textbf{\color{accentcopper}아레스:} #1\par}
\newcommand{\sayalbi}[1]{\textbf{\color{skyblue}알비:} #1\par}

%% 미션 메타 정보 헬퍼
\newcommand{\missionmeta}[2]{%
  \par\noindent
  {\sffamily\small\color{neutralgray}\textbf{사용 부품:} #1 \quad|\quad \textbf{학습 목표 분류:} #2}\par\vspace{0.4em}%
}
